% GENERATED FILE. Edit canonical lessons, then run npm run generate:books.
% canonical-source-hash: fnv1a64:87ee03966928d46a
% canonical-lessons: PA-C62-guauna, PA-C62-pahinna, PA-C62-nahauna, PA-C62-chhaddna, PA-C62-udikna

\chapter{To lose, To wear, To bathe, To leave, To wait for}
\label{ch:pa-to-lose-to-wear-to-bathe-to-leave-to-wait-for}
\hlchaptermodality{62}

\begin{chapteropening}
\textbf{By the end of this chapter:} \emph{I can say to lose, to wear, to bathe, to leave, to let go and to wait for.}

Complete the last lesson of chapter 62: I can say to lose, to wear, to bathe, to leave, to let go and to wait for.
\end{chapteropening}

\section[\texorpdfstring{\pa{ਗੁਆਉਣਾ} (guāuṇā)}{ਗੁਆਉਣਾ (guāuṇā)}]{\pa{ਗੁਆਉਣਾ} (guāuṇā) — to lose}

\label{lesson:PA-C62-guauna}



\begin{quote}
Before the new one: say the Punjabi for to stop, then the Punjabi for to tell.
\end{quote}

\subsection*{You'll want to know: \pa{ਗੁਆਉਣਾ}}
\textbf{\pa{ਗੁਆਉਣਾ}} — \emph{guāuṇā} — \textquotedblleft{}to lose\textquotedblright{}.

\begin{cousinweb}
\textbf{\pa{ਗੁਆਉਣਾ}} (\emph{guāuṇā}) is inherited. CDIAL lists it, in the spelling \emph{gãvāuṇā}, \textquotedblleft{}to lose, damage\textquotedblright{}, under Sanskrit \emph{gamayati}, \textquotedblleft{}causes to go\textquotedblright{} (CDIAL 4028): to lose a thing is to make it go.
\end{cousinweb}

\begin{grammarlens}[title={What you've built}]
One more word for this chapter.
\end{grammarlens}

\subsection*{Your turn}
\begin{itemize}
\raggedright
  \item \emph{Say it:} \emph{guāuṇā}
  \item \emph{Say it:} \emph{guāuṇā}, once more
  \item \emph{Say it:} say \emph{guāuṇā}
  \item \emph{Recall:} say the Punjabi for to stop, then the Punjabi for to tell, then say \emph{guāuṇā} again
\end{itemize}

\begin{tcolorbox}[breakable,colback=teal!4,colframe=teal!35!black,title={Before you move on}]
What does \pa{ਗੁਆਉਣਾ} mean? (\textquotedblleft{}To lose\textquotedblright{}.) Say it once more.
\end{tcolorbox}
\section[\texorpdfstring{\pa{ਪਹਿਨਣਾ} (pahinṇā)}{ਪਹਿਨਣਾ (pahinṇā)}]{\pa{ਪਹਿਨਣਾ} (pahinṇā) — to wear}

\label{lesson:PA-C62-pahinna}



\begin{quote}
Before the new one: say the Punjabi for to tell, then the Punjabi for to lose.
\end{quote}

\subsection*{You'll want to know: \pa{ਪਹਿਨਣਾ}}
\textbf{\pa{ਪਹਿਨਣਾ}} — \emph{pahinṇā} — \textquotedblleft{}to wear\textquotedblright{}.

\begin{cousinweb}
\textbf{\pa{ਪਹਿਨਣਾ}} (\emph{pahinṇā}) is a tadbhava, a word inherited and reshaped by sound change. It comes from \emph{*pinahati}, \textquotedblleft{}ties on\textquotedblright{}, a form CDIAL reconstructs from its descendants (CDIAL 8198). It grew out of Sanskrit \emph{pinahya}, \textquotedblleft{}having dressed\textquotedblright{}.
\end{cousinweb}

\begin{grammarlens}[title={What you've built}]
One more word for this chapter.
\end{grammarlens}

\subsection*{Your turn}
\begin{itemize}
\raggedright
  \item \emph{Say it:} \emph{pahinṇā}
  \item \emph{Say it:} \emph{pahinṇā}, once more
  \item \emph{Say it:} say \emph{pahinṇā}
  \item \emph{Recall:} say the Punjabi for to tell, then the Punjabi for to lose, then say \emph{pahinṇā} again
\end{itemize}

\begin{tcolorbox}[breakable,colback=teal!4,colframe=teal!35!black,title={Before you move on}]
What does \pa{ਪਹਿਨਣਾ} mean? (\textquotedblleft{}To wear\textquotedblright{}.) Say it once more.
\end{tcolorbox}
\section[\texorpdfstring{\pa{ਨਹਾਉਣਾ} (nahāuṇā)}{ਨਹਾਉਣਾ (nahāuṇā)}]{\pa{ਨਹਾਉਣਾ} (nahāuṇā) — to bathe}

\label{lesson:PA-C62-nahauna}



\begin{quote}
Before the new one: say the Punjabi for to lose, then the Punjabi for to wear.
\end{quote}

\subsection*{You'll want to know: \pa{ਨਹਾਉਣਾ}}
\textbf{\pa{ਨਹਾਉਣਾ}} — \emph{nahāuṇā} — \textquotedblleft{}to bathe\textquotedblright{}.

\begin{cousinweb}
\textbf{\pa{ਨਹਾਉਣਾ}} (\emph{nahāuṇā}) is inherited. CDIAL lists it, as \emph{nhāuṇā}, under Sanskrit \emph{snāpayati}, \textquotedblleft{}causes to bathe, washes\textquotedblright{} (CDIAL 13791).
\end{cousinweb}

\begin{grammarlens}[title={What you've built}]
One more word for this chapter.
\end{grammarlens}

\subsection*{Your turn}
\begin{itemize}
\raggedright
  \item \emph{Say it:} \emph{nahāuṇā}
  \item \emph{Say it:} \emph{nahāuṇā}, once more
  \item \emph{Say it:} say \emph{nahāuṇā}
  \item \emph{Recall:} say the Punjabi for to lose, then the Punjabi for to wear, then say \emph{nahāuṇā} again
\end{itemize}

\begin{tcolorbox}[breakable,colback=teal!4,colframe=teal!35!black,title={Before you move on}]
What does \pa{ਨਹਾਉਣਾ} mean? (\textquotedblleft{}To bathe\textquotedblright{}.) Say it once more.
\end{tcolorbox}
\section[\texorpdfstring{\pa{ਛੱਡਣਾ} (chhaḍḍṇā)}{ਛੱਡਣਾ (chhaḍḍṇā)}]{\pa{ਛੱਡਣਾ} (chhaḍḍṇā) — to leave, to let go}

\label{lesson:PA-C62-chhaddna}



\begin{quote}
Before the new one: say the Punjabi for to wear, then the Punjabi for to bathe.
\end{quote}

\subsection*{You'll want to know: \pa{ਛੱਡਣਾ}}
\textbf{\pa{ਛੱਡਣਾ}} — \emph{chhaḍḍṇā} — \textquotedblleft{}to leave, to let go\textquotedblright{}.

\begin{cousinweb}
\textbf{\pa{ਛੱਡਣਾ}} (\emph{chhaḍḍṇā}) is inherited. CDIAL lists it under Sanskrit \emph{chardayati}, \textquotedblleft{}causes to flow over; vomits\textquotedblright{}, a verb whose basic meaning it gives as \textquotedblleft{}eject\textquotedblright{} (CDIAL 4998).
\end{cousinweb}

\begin{grammarlens}[title={What you've built}]
One more word for this chapter.
\end{grammarlens}

\subsection*{Your turn}
\begin{itemize}
\raggedright
  \item \emph{Say it:} \emph{chhaḍḍṇā}
  \item \emph{Say it:} \emph{chhaḍḍṇā}, once more
  \item \emph{Say it:} say \emph{chhaḍḍṇā}
  \item \emph{Recall:} say the Punjabi for to wear, then the Punjabi for to bathe, then say \emph{chhaḍḍṇā} again
\end{itemize}

\begin{tcolorbox}[breakable,colback=teal!4,colframe=teal!35!black,title={Before you move on}]
What does \pa{ਛੱਡਣਾ} mean? (\textquotedblleft{}To leave, to let go\textquotedblright{}.) Say it once more.
\end{tcolorbox}
\section[\texorpdfstring{\pa{ਉਡੀਕਣਾ} (uḍīkṇā)}{ਉਡੀਕਣਾ (uḍīkṇā)}]{\pa{ਉਡੀਕਣਾ} (uḍīkṇā) — to wait for}

\label{lesson:PA-C62-udikna}



\begin{quote}
Before the new one: say the Punjabi for to bathe, then the Punjabi for to leave.
\end{quote}

\subsection*{You'll want to know: \pa{ਉਡੀਕਣਾ}}
\textbf{\pa{ਉਡੀਕਣਾ}} — \emph{uḍīkṇā} — \textquotedblleft{}to wait for\textquotedblright{}.

\begin{cousinweb}
\textbf{\pa{ਉਡੀਕਣਾ}} (\emph{uḍīkṇā}) is a tadbhava, a word inherited and reshaped by sound change. It comes from \emph{*uddekṣati}, \textquotedblleft{}looks for, expects\textquotedblright{}, a form CDIAL reconstructs from its descendants (CDIAL 1997).
\end{cousinweb}

\begin{grammarlens}[title={What you've built}]
One more word for this chapter.
\end{grammarlens}

\subsection*{Your turn}
\begin{itemize}
\raggedright
  \item \emph{Say it:} \emph{uḍīkṇā}
  \item \emph{Say it:} \emph{uḍīkṇā}, once more
  \item \emph{Say it:} say \emph{uḍīkṇā}
  \item \emph{Recall:} say the Punjabi for to bathe, then the Punjabi for to leave, then say \emph{uḍīkṇā} again
\end{itemize}

\begin{tcolorbox}[breakable,colback=teal!4,colframe=teal!35!black,title={Before you move on}]
What does \pa{ਉਡੀਕਣਾ} mean? (\textquotedblleft{}To wait for\textquotedblright{}.) Say it once more.
\end{tcolorbox}
